\setchapterimage{figure/Chapter11封面图.jpg}
\setchapterintro{本章从坐标图定义流形，再构造切空间、切丛和向量场。在坐标图中可应用前面建立的微分方程局部理论；紧性则使局部解拼接成全局相流。}
\chapter{流形上的微分方程}\label{chap:manifolds}

第~\ref{chap:arnold-geometry}~章中的坐标变换发生在 Euclidean 开集之间；环面运动又表明，一个相空间未必能由单张 Euclidean 坐标图覆盖。本章先建立换图、切空间和切丛，再把已有局部方程拼接为流形上的相流。

\section{坐标图与微分结构}
\begin{definition}
集合 $M$ 的一张 $n$ 维坐标图是双射 $\phi:W\to U$，其中 $W\subset M$、
$U\subset\R^n$ 为开集。
两图相容，是指重叠部分在各图中的像为开集，且过渡映射
$\phi_j\phi_i^{-1}$ 为光滑微分同胚。
相容且覆盖 $M$ 的坐标图族称图册。
两个图册若合并后仍相容，称为等价；等价类称光滑结构。
带光滑结构的集合，在下述分离与可数条件下称光滑流形。
\end{definition}
\begin{definition}
流形子集 $G$ 称开集，若每图中的 $\phi(G\cap W)$ 均开。
要求任意两不同点有不相交开邻域，称 Hausdorff 条件。
要求拓扑有可数个开集组成的基，称第二可数条件；这里“基”指任一开集都能写成这些开集的并。
对于局部 Euclidean 空间，这等价于存在可数坐标图册：每张图内的有理球给出可数基；
反向则用可数基选取包含在坐标域内的基元素，得到可数坐标子覆盖。
本章的流形均满足这两个条件，且无边界。
\end{definition}
由过渡映射相容性可知开集定义与图册选择无关。
若仅要求过渡映射为同胚，得到拓扑流形；要求 $C^r$ 或解析则得到相应类别。
定向图册要求换图映射的 Jacobi 行列式为正。


\begin{definition}[定向、闭流形与带边界流形]\label{def:orientation-boundary}
在 $n$ 维实向量空间中，两组有序基若换基行列式为正，就规定它们有相同定向。
在每个坐标域选择定向，且所有重叠处换图行列式为正，称给流形选定了定向；
能够这样选择时称可定向。
紧且无边界的流形称闭流形，这里的“闭”不只是说它在某个环境空间中是闭集。
允许坐标域是半空间 $\{x_n\geq0\}$ 中的相对开集，就得到带边界流形；
局部对应 $x_n=0$ 的点组成边界 $\partial M$。
边界的定向约定为：先写向外法向量，再写边界的有序切向基，应得到内部的正向基。
\end{definition}

\begin{example}[球面的两图]
单位球面 $S^n\subset\R^{n+1}$ 的南北极立体投影给出两图。
其过渡映射在非零点为 $x\mapsto x/|x|^2$，自身是逆，故光滑。
这两图覆盖整个球面，定义其光滑结构。
\end{example}
\begin{example}[射影空间]
$\mathbb{RP}^n$ 是 $\R^{n+1}$ 中过原点的直线的集合。
齐次坐标 $[x_0:\cdots:x_n]$ 把非零数倍坐标视为同一点。
在 $x_i\ne0$ 的区域，取其余 $n$ 个比值 $x_j/x_i$ 为坐标。
过渡关系均为分母非零的有理函数，所以给出 $n+1$ 张相容图。
\end{example}
环面用每个角坐标的两个开区间作图，其乘积给出有限图册。
这一构造不要求先把环面画在三维空间中。


\begin{definition}[覆盖映射]\label{def:covering-map}
连续满射 $\pi:\widetilde M\to M$ 称覆盖映射，若每个 $x\in M$ 有开邻域 $U$，
使 $\pi^{-1}(U)$ 是两两不交开集的并，且 $\pi$ 在每个这样的开集上都同胚到 $U$。
若每个点恰有 $k$ 个原像，称 $k$ 重覆盖。
例如 $\R\to\R/2\pi\Z$，$t\mapsto[t]$ 是覆盖；$S^n\to\mathbb{RP}^n$，$x\mapsto[x]$ 是双覆盖。
在半径足够小的球面片上，片与其对径片不交，各片都能唯一选取单位向量代表，故满足定义。
\end{definition}

\begin{definition}
子集 $K\subset M$ 紧，是指它的每个开覆盖均有有限子覆盖。
流形连通，是指它不能分成两个不相交非空开集。
连通分支的维数由任一坐标图的 Euclidean 维数定义。
\end{definition}
\begin{proposition}
若 $K$ 是有限个子集 $K_i$ 的并，各 $K_i$ 包含在一张图内且图像紧，则 $K$ 紧。
连通流形上的所有坐标图维数相同。
\end{proposition}
\begin{proof}
对任意开覆盖，利用坐标映射将各开集表示在相应的坐标域中，再由 $\phi_i(K_i)$ 的紧性选有限子覆盖，
再对有限个 $i$ 合并，得到 $K$ 的有限子覆盖。
换图映射的微分是两个坐标空间间的线性同构，故相交图具有相同维数。
固定一点，从它可由有限条相交坐标邻域链到达的点集是开集，其补集也开。
连通性使它为整个流形，所以所有图可沿链比较维数。
\end{proof}
球面以有限个有界小坐标片的紧闭子片覆盖，因而紧；射影空间亦可由单位球面
经连续映射 $x\mapsto[x]$ 得到，故紧。
局部图为局部道路连通空间，所以上述连通性也等价于有限图链连通和道路连通。


图~\ref{fig:ode-tikz-stereographic} 画出球面的一张坐标图；北极自身不在这张图的定义域内。
\begin{center}
\begin{minipage}{0.96\linewidth}
\centering
% TikZ mathematical illustration; compiled with the main document.
\begin{tikzpicture}[>=Stealth,font=\small,scale=2.15]
 \fill[GreenPale,opacity=.3] (0,0) circle(1);
 \draw[GreenMain,thick] (0,0) circle(1);
 \draw[->,gray] (-1.45,0)--(2.55,0) node[right]{投影平面};
 \draw[gray,dashed] (0,-1.15)--(0,1.25);
 \draw[GoldMain,thick] (0,1)--(2,0);
 \draw[GoldMain,->,thick] (.95,.525)--(1.35,.325);
 \fill (0,1) circle(1pt) node[above left]{$N$};
 \fill[GreenMain] (.8,.6) circle(1pt) node[above right]{$P$};
 \fill[GoldMain] (2,0) circle(1pt) node[below]{$q$};
 \node[below left] at (0,0){$0$};
\end{tikzpicture}
\captionof{figure}{立体投影的二维截面：从北极 $N$ 出发的直线将球面点 $P$ 映到平面点 $q$。}
\label{fig:ode-tikz-stereographic}
\end{minipage}
\end{center}

\section{光滑映射、子流形与例子}
\begin{definition}
$f:M\to N$ 光滑，是指其所有局部坐标表示 $\psi f\phi^{-1}$ 光滑。
光滑可逆且逆也光滑时称微分同胚。
光滑曲线是区间到流形的光滑映射，可有自交或零速度。
子集 $V\subset M$ 称 $k$ 维子流形，若在其每点附近有图使它成为
$\R^k\times\{0\}$ 与图域的交。
\end{definition}
\begin{proposition}[正则方程定义子流形]
若 $F:U\subset\R^N\to\R^{N-k}$ 光滑，且在 $F^{-1}(0)$ 上微分满秩，
则 $F^{-1}(0)$ 是 $k$ 维子流形。
\end{proposition}
\begin{proof}
在一点选取一个非零的 $(N-k)$ 阶子式，把其对应变量与其余 $k$ 个变量分开。
隐函数定理把前者写成后者的光滑函数。
以“前者减去该函数”和后者作坐标，零集恰成为坐标平面。
\end{proof}
下文将实正交群及特殊正交群分别记为 $O(n)$ 和 $SO(n)$。
\begin{example}[正交群]
把矩阵视为 $\R^{n^2}$ 的点，令 $F(Q)=Q^TQ-I$，其值在对称矩阵空间。
在正交矩阵处，$DF_Q(H)=Q^TH+H^TQ$。
任给对称矩阵 $S$，取 $H=QS/2$ 即得到 $S$，故满秩。
所以 $O(n)$ 是维数 $n(n-1)/2$ 的流形。
行列式的符号区分两个连通分支：$SO(n)$ 与一个反射乘 $SO(n)$。
用二维坐标平面旋转逐次消去一列的非对角元，可将任意 $SO(n)$ 元素连续连到恒等阵；
最后剩余对角的成对负号可用角度从零到 $\pi$ 的平面旋转消去。
因此这两个集合确实分别连通。
\end{example}
同样，$GL(n,\R)$ 也有两个分支。对可逆矩阵作极分解
$A=QP$，$P=(A^TA)^{1/2}>0$，路径 $(1-t)P+tI$ 始终正定，
把它变形到 $Q$，再用上面对正交群的分析即可。
取 $A^TA$ 的正定平方根，并令 $Q=AP^{-1}$，则
$Q^TQ=P^{-1}A^TAP^{-1}=I$，所以这也给出了所用极分解的具体构造。


\begin{definition}[浸入、嵌入与把手]\label{def:immersion-embedding}
光滑映射 $f:M\to N$ 的每点微分单射时称浸入；若它还单射，且作为 $M\to f(M)$ 的映射是到子空间拓扑的同胚，称光滑嵌入。
因此嵌入不仅要求参数化没有零速度，还要求没有不同点重合或全局拓扑不相容。
从两个曲面各去掉一个开圆盘，再沿边界圆粘合，得到它们的连通和。
球面与 $g$ 个环面的连通和称为带 $g$ 个把手的可定向曲面；$g$ 称亏格。
\end{definition}

\paragraph{拓扑预备结论。}
连通一维无边界流形微分同胚于圆或实线；闭连通可定向曲面微分同胚于带若干把手的球面。
Whitney 嵌入定理保证每个光滑 $n$ 维流形能嵌入充分高维的 Euclidean 空间。
这些分类与嵌入结果在本章中作为拓扑预备定理使用。
流形同胚不必微分同胚；同一拓扑球面上可能存在不同的光滑结构。

\section{切空间与切丛}
\begin{definition}
过 $x\in M$ 的两条光滑曲线，若在一张图中的初速度相同，称在 $x$ 相切。
等价类称 $x$ 处切向量，所有等价类组成 $T_xM$。
其线性运算由坐标中的速度向量定义。
\end{definition}
\begin{proposition}
相切与线性结构不依赖图的选择；$\dim T_xM=\dim M$。
光滑映射 $f$ 诱导线性映射
\[
 T_xf:T_xM\to T_{f(x)}N,
\]
且 $T(gf)=Tg\circ Tf$。
\end{proposition}
\begin{proof}
另一张图中的速度由换图映射的 Jacobi 矩阵作用于原速度得到，该矩阵可逆。
所以速度相等、加法和数乘在换图后保持。
每个坐标速度均可由一条小直线的逆图像实现，故切空间同构于 $\R^n$。
令 
\[
 T_xf[\gamma]=[f\gamma],
\]
链式法则同时证明良定义、线性和复合公式。
\end{proof}
\begin{definition}
不交并 
\[
 TM=\coprod_{x\in M}T_xM
\]
 称切丛的总空间。
投影 $p:TM\to M$ 把切向量送到附着点，$M$ 称底空间，$p^{-1}(x)$ 称纤维。
零截面 $i:M\to TM$ 将 $x$ 送到 $T_xM$ 的零向量。
\end{definition}
\begin{theorem}
$TM$ 有自然的 $2n$ 维光滑流形结构；其图坐标为 $(x,v)$，换图为
\[
 (x,v)\longmapsto(g(x),Dg(x)v).
\]
$p,i$ 及 $Tf:TM\to TN$ 都光滑，且 $pi=\operatorname{id}_M$。
\end{theorem}
\begin{proof}
换图公式来自速度的链式法则，光滑且逆由 $g^{-1}$ 给出。
复合过渡关系也由链式法则满足，所以给出图册。
在图中，投影为 $(x,v)\mapsto x$，零截面为 $x\mapsto(x,0)$，
切映射为 $(x,v)\mapsto(f(x),Df(x)v)$，因此均光滑，复合式也直接成立。
\end{proof}
对于 $C^r$ 流形，切丛的换图映射一般仅属于 $C^{r-1}$，因为其表达式含原换图映射的一阶导数。本章讨论 $C^\infty$ 流形。


\begin{definition}[向量丛、截面与标架]\label{def:vector-bundle}
秩 $k$ 的光滑向量丛包括总空间 $E$、底空间 $M$ 和光滑投影 $p:E\to M$，
每个纤维 $E_x=p^{-1}(x)$ 是 $k$ 维向量空间；在每点附近有
$p^{-1}(U)\simeq U\times\R^k$ 的光滑对应，且在每条纤维上是线性同构。
这样的对应称局部平凡化。光滑映射 $s:M\to E$ 满足 $p\circ s=\operatorname{id}$ 时称截面。
两个同底空间向量丛之间逐纤维线性的微分同胚称丛同构。
一组在每条纤维组成基的光滑截面称标架；只在一个开集上选取时称局部标架。
切丛的局部平凡化就是前面构造的 $(x,v)$ 坐标，因而满足这些条件。
\end{definition}
\begin{definition}[度量与梯度]\label{def:riemannian-gradient}
Riemann 度量是在各切空间上光滑选取的正定内积 $g_x$。
在局部坐标中它由正定对称矩阵 $(g_{ij}(x))$ 表示，
$|v|_g^2=\sum_{i,j}g_{ij}v_iv_j$。
函数 $f$ 的梯度由
\[
 g_x(\operatorname{grad}f(x),v)=df_x(v)\quad(v\in T_xM)
\]
唯一确定，坐标分量是 $(g_{ij})^{-1}Df^T$。
对于 Euclidean 子流形，使用环境内积在切空间上的限制。
单位切向量丛为 $T^1M=\{(x,v)\in TM:|v|_g=1\}$。
\end{definition}

\begin{definition}
若存在逐纤维线性的丛同构 $TM\simeq M\times\R^n$，称 $M$ 可平行化。
等价地，可在全流形选出 $n$ 个处处线性无关的光滑向量场。
\end{definition}
\begin{proof}[等价性的证明]
丛同构把 $\R^n$ 的标准基拉回，得到全局标架。
反之，给定标架，把 $(x,c_1,\ldots,c_n)$ 送到基向量的线性组合。
在局部图中这是光滑可逆矩阵作用，逆矩阵光滑，所以得到丛同构。
\end{proof}
Euclidean 开集与环面可平行化，后者取各角方向的速度场。

\begin{definition}[Möbius 带]
将 $[0,1]\times(-1,1)$ 中的 $(0,s)$ 与 $(1,-s)$ 识别，取商拓扑及由局部矩形诱导的光滑坐标，
所得曲面称开 Möbius 带。区间的中线拼成圆，横向坐标绕一周变号。
\end{definition}

Möbius 带不可平行化：沿中心圆走一周，横向基方向反转，换图映射的 Jacobi 行列式为负，
所以它不可定向；而全局标架会给出连续定向，矛盾。
球面 $S^2$ 不存在无零切向量场，将在指标和定理后证明。
对于任意维球面，可平行化性需要另行讨论。这里仅使用上述例子及后面由指标和定理得到的 $S^2$ 结论。

\section{几个相互微分同胚的三维流形}
\begin{definition}[四元数的共轭与范数]
四元数的标准基记为 $1,i,j,k$。对 $q=a+bi+cj+dk$，定义其共轭为
$\overline q=a-bi-cj-dk$。
范数由 $|q|^2=q\overline q$ 给出，纯虚四元数指实部为零的元素。
结合律及 $\overline{pq}=\overline q\,\overline p$ 给出
$|pq|^2=pq\overline q\,\overline p=|p|^2|q|^2$。
\end{definition}
\begin{proposition}
下列流形微分同胚：$SO(3)$、单位切向量丛 $T^1S^2$、$\mathbb{RP}^3$、
空间中固定一点的刚体的姿态空间，以及
\[
 \{z\in\C^3:z_1^2+z_2^2+z_3^2=0,\quad |z|^2=2\}.
\]
\end{proposition}
\begin{proof}
给 $(x,v)\in T^1S^2$，$|x|=|v|=1,x\perp v$，矩阵 $(x,v,x\times v)$ 属于 $SO(3)$。
取前两列是逆映射，故是微分同胚。刚体姿态由一组正向正交标架确定，
也与 $SO(3)$ 对应。
写 $z=x+iy$，复方程给出 $|x|^2=|y|^2,x\cdot y=0$，
范数条件给出二者均为 1，因而与同一个单位切向量模型对应。

最后用单位四元数 $q\in S^3$ 作用于纯虚四元数 $v\in\R^3$：
$v\mapsto qvq^{-1}$。乘法保持范数，所以得到旋转；
绕单位轴 $u$ 旋转角 $\theta$ 由 $q=\cos(\theta/2)+u\sin(\theta/2)$ 实现。
核恰为 $\{1,-1\}$，故诱导 $S^3/\{\pm1\}=\mathbb{RP}^3$ 到 $SO(3)$ 的双射。
其在单位元的微分为 $v\mapsto2u\times v$，从纯虚四元数到反对称矩阵空间是同构，
所以是局部微分同胚；结合双射得到全局微分同胚。
\end{proof}

\section{向量场与紧支撑相流}
\begin{definition}
流形上的向量场是光滑截面 $v:M\to TM$，满足 $pv=\operatorname{id}_M$。
积分曲线满足 $\gamma'(t)=v(\gamma(t))$。
若 $v$ 在某紧集外为零，称为紧支撑向量场；若其所有积分曲线均双向全局存在，称为完备。
\end{definition}
\begin{theorem}[紧支撑完备性]\label{thm:arnold-compact-flow}
紧支撑光滑向量场完备，并生成全局单参数微分同胚群。
特别地，紧流形上的每个光滑向量场都完备。
\end{theorem}
\begin{proof}
每点在一张坐标图内应用前文的存在唯一性与光滑依赖定理，
得到一个邻域及所有初值共用的正存在时间。
这里 $K=\operatorname{supp}v$，支撑按定义~\ref{def:support-test} 在流形拓扑中取闭包。用有限个这样的邻域覆盖紧集 $K$，取这些时间的一个正下界 $\delta$。
在 $K$ 外向量场为零，积分曲线恒定。
所以每个初值都至少有共同时间 $(-\delta,\delta)$ 的解。

在重叠坐标图里，局部解由唯一性一致，故拼接成流形上的解。
Hausdorff 条件保证“相同到某时刻”的闭性，排除在分叉的非 Hausdorff 空间中
局部唯一却整体分离的情况。
把小时间解反复复合，任意有限时间都可分成有限个长度小于 $\delta$ 的步长，
所以得到所有 $t\in\R$ 的解。
唯一性给出群性质，$\varphi^{-t}$ 为逆；每个有限复合由局部光滑映射组成，故光滑。
\end{proof}
\begin{corollary}
若 $M$ 紧，且 $v(t,x)$ 在 $\R\times M$ 光滑，则非自治方程的每个解也双向全局存在。
\end{corollary}
\begin{proof}
若最大存在区间有有限端点 $T$，在紧集 $[T-1,T+1]\times M$ 上用有限图覆盖，
得到共同正存在时间。取接近 $T$ 的一个已有时刻，其局部解可以跨越 $T$，矛盾。
正反向端点的证明相同。
\end{proof}

\section{线性方程的射影紧化}
\begin{proposition}
线性向量场 $x'=A(t)x$ 从 $\R^n$ 延伸为 $\mathbb{RP}^n$ 上的光滑含时向量场，
且无穷远超平面不变。
\end{proposition}
\begin{proof}
在无穷远附近取坐标 $y_1=1/x_1,y_k=x_k/x_1$（$k>1$）。
记 $B(y)=a_{11}+\sum_{j>1}a_{1j}y_j$，逐项求导得
\[
 y_1'=-y_1B(y),\qquad
 y_k'=a_{k1}+\sum_{j>1}a_{kj}y_j-y_kB(y).
\]
右端是 $y$ 的多项式，可以光滑延伸到 $y_1=0$。
那里第一分量为零，所以无穷远超平面不变。
其余图同理，重叠区域因与原场一致而相容。
\end{proof}
射影空间紧，故其含时解全局存在。
从仿射部分出发的轨道不可能在有限时间到达无穷远超平面：若到达，
由反向唯一性，它原本就应在这个不变超平面内。
因此原线性方程的解全局存在，这也是前面解的延拓定理的结论。
